\documentclass[11pt]{article}

% --------------------------------------------------
% Packages
% --------------------------------------------------
\usepackage[T1]{fontenc}
\usepackage{lmodern}
\usepackage{geometry}
\usepackage{setspace}
\usepackage{amsmath, amssymb, amsthm}
\usepackage{csquotes}
\usepackage{hyperref}

\geometry{margin=1in}
\setstretch{1.15}

% --------------------------------------------------
% Theorem environments
% --------------------------------------------------
\newtheorem{definition}{Definition}
\newtheorem{proposition}{Proposition}
\newtheorem{theorem}{Theorem}

% --------------------------------------------------
% Title information
% --------------------------------------------------
\title{Neural Commit Semantics and the Algebraic Structure of Language}
\author{Flyxion}
\date{December 2025}

\begin{document}

\maketitle

\begin{abstract}
Human language exhibits algebraic properties—most notably commutativity with respect to linear order and nonassociativity under hierarchical composition—that resist explanation in purely statistical or state-based models. Recent neurocomputational theories, including Murphy’s ROSE framework, argue that these properties are enforced by neural control dynamics rather than symbolic rules. In this paper, we provide a formal account of this claim by introducing \emph{neural commit semantics}: an event-first architecture in which irreversible commit operations enforce algebraic invariants over replayable histories. We show that commutativity arises naturally from order-insensitive merge under canonical collapse, while nonassociativity is enforced by commit-induced irreversibility. We reinterpret the ROSE architecture as a perceptual control system operating over an authoritative event history, and we argue that contemporary large language models lack the architectural capacity to enforce these invariants under counterfactual intervention. The result is a unified account of linguistic structure, neural dynamics, and explanation itself, grounded in counterfactual control rather than prediction.
\end{abstract}

\section{Introduction}

Human language combines unbounded generativity with strict structural constraint. Speakers effortlessly construct and interpret expressions they have never encountered before, yet those expressions obey precise algebraic regularities. Among the most striking of these are two properties emphasized throughout generative linguistics: syntactic composition is commutative with respect to linear order, yet nonassociative with respect to hierarchical grouping (Chomsky 1995).

Recent advances in large language models have sharpened, rather than resolved, the explanatory challenge posed by these properties. Although such models achieve impressive predictive performance and correlate with neural measurements, they systematically fail to enforce higher-order structural constraints under counterfactual manipulation. This raises a central question: what kind of architecture is required not merely to predict linguistic behavior, but to explain the algebraic structure that defines it?

A growing body of theoretical work suggests that explanation must appeal to control rather than prediction. In particular, Murphy’s ROSE (Representation, Operation, Structure, Encoding) framework proposes that linguistic structure is enforced by neural coordination dynamics operating across multiple timescales, rather than by stored symbolic representations (Murphy 2023). On this view, algebraic properties such as headedness and nonassociativity emerge from the timing and irreversibility of neural events.

In this paper, we formalize that intuition. We argue that linguistic algebra is enforced by \emph{commit semantics}: irreversible control operations that stabilize speculative composition into authoritative structure. We adopt an event-first ontology in which semantic structure arises from replayable histories rather than mutable states, and we show that commutativity and nonassociativity follow necessarily from architectural constraints on such histories. Explanation, on this view, is identified not with prediction or correlation, but with the capacity to evaluate counterfactual interventions.

\section{Why Prediction Is Not Explanation}

Prediction and explanation are often conflated in contemporary computational models of cognition. Systems that accurately predict behavior or neural activity are frequently taken to have explained the underlying phenomenon. However, work in causal inference has made clear that prediction alone is insufficient for explanation (Pearl 2009; Pearl and Mackenzie 2018).

Explanation requires counterfactual competence: the ability to determine what would happen under interventions that alter the underlying causal structure. A model that merely tracks correlations may succeed within the training distribution while failing catastrophically when asked whether a particular structural manipulation is admissible.

Language provides a particularly sharp illustration of this distinction. Many linguistic expressions are structurally ambiguous: a single surface string may correspond to multiple distinct hierarchical organizations, each internally well-formed. The explanatory task is not to predict the most likely interpretation, but to account for why certain structures are permitted at all, while others are categorically excluded.

Statistical sequence models collapse this distinction. Alternative parses are represented as probability mass over continuations, not as distinct admissible histories. As a result, such models conflate structural ambiguity with uncertainty, and they lack the means to rule out structurally impossible configurations.

\section{An Event-First Criterion for Explanation}

To make the notion of explanation precise, we adopt an event-first ontology. Rather than treating cognition as a sequence of state transitions, we model it as the construction of replayable histories composed of discrete events.

\begin{definition}[Event History]
An event history is a partially ordered set of events together with dependency relations specifying how later events are licensed by earlier ones.
\end{definition}

Explanation, in this framework, consists in specifying which histories are admissible and which are not. A system explains a domain if it can determine whether a proposed intervention yields a valid history.

\begin{theorem}[Counterfactual Criterion for Explanation]
A model explains a structured domain if and only if it can evaluate the admissibility of counterfactual interventions over event histories.
\end{theorem}

This criterion distinguishes models that merely approximate surface regularities from those that enforce underlying structural constraints.

\section{Commit Semantics and Linguistic Algebra}

We now show how the algebraic properties of language follow from commit semantics in an event-based architecture.

\subsection{Merge as Event Union}

We treat basic linguistic composition as the union of event-regions rather than the application of a symbolic operator.

\begin{definition}[Event-Based Merge]
Given two event-regions $x$ and $y$, merge is defined as region union:
\[
\mathsf{Merge}(x,y) = x \cup y .
\]
\end{definition}

Because union is symmetric, merge is commutative by construction. At this stage, no order or headedness is fixed.

\subsection{Commit as Irreversible Collapse}

Structure becomes stable only through commit.

\begin{definition}[Commit]
A commit operation applies a canonical collapse operator $D$ to a composite region, yielding an indivisible unit whose internal structure is no longer accessible.
\end{definition}

Commit is an irreversible control event, not a representational update.

\begin{proposition}[Commit-Enforced Nonassociativity]
Once a composite region has been committed, reassociation of its internal constituents is structurally impossible.
\end{proposition}

\begin{proof}
Consider $D(D(x \cup y) \cup z)$ and $D(x \cup D(y \cup z))$. Because commit renders internal structure inaccessible, there exists no sequence of admissible operations transforming one into the other without violating irreversibility. Hence the two are distinct.
\end{proof}

Headedness arises at commit time via dominance relations that are invariant under permutation, yielding commutativity without stipulation.

\section{ROSE as Neural Perceptual Control}

We now reinterpret the ROSE architecture as a perceptual control system rather than a processing pipeline.

The R component provides access to available event indices; O generates speculative merge proposals; S specifies structural invariants to be maintained; and E coordinates stabilization across distributed neural systems. None of these components stores structure. Instead, they regulate when structure becomes fixed.

Commit events act as the pivot of this control loop. Neurally, these correspond to brief coordination bursts—such as beta-mediated events—that terminate speculation, stabilize structure, and prepare the system for the next compositional cycle (Murphy 2023).

On this view, syntax is not localized to a region, but realized as the maintenance of invariants through coordinated neural dynamics.

\section{Why Large Language Models Fail}

Large language models optimize predictive accuracy over token sequences. All internal representations remain simultaneously accessible, and no operation renders a subset of structure irrevocable.

\begin{proposition}
Large language models cannot enforce nonassociativity under counterfactual intervention.
\end{proposition}

\begin{proof}
Transformer architectures maintain access to all token representations at every layer via attention. No mechanism collapses a subset into an indivisible unit whose internal structure cannot be revisited. Hence reassociation is always possible.
\end{proof}

Apparent commutativity in such models reflects properties of training data, not architectural enforcement. Scaling improves approximation of surface distributions but does not introduce commit semantics or authoritative histories.

\section{Discussion}

The account developed here reframes longstanding debates about language, cognition, and computation. It explains how abstract algebraic properties can be neurally enforced without symbolic encoding, why oscillatory dynamics are central, and why prediction alone is insufficient for explanation.

The framework is compatible with causal approaches to cognition (Pearl 2009), with control-theoretic perspectives, and with multi-level analyses in the spirit of Marr (1982). It also clarifies why certain hybrid neurocomputational models succeed where purely statistical ones fail.

\section{Conclusion}

We have argued that the algebraic structure of language is enforced by neural commit semantics operating over event histories. Commutativity and nonassociativity emerge from architectural constraints, not from symbolic stipulation or statistical learning. Explanation, on this view, consists in counterfactual control rather than prediction. Language marks a phase transition in cognition not because it introduces new representations, but because it introduces irreversible commitment.


\begin{thebibliography}{99}

\bibitem{murphy2023rose}
Murphy, E. (2023).
\newblock \emph{ROSE: A Neurocomputational Architecture for Syntax}.
\newblock arXiv preprint arXiv:2303.08877.

\bibitem{chomsky1995minimalist}
Chomsky, N. (1995).
\newblock \emph{The Minimalist Program}.
\newblock MIT Press.

\bibitem{pearl2009causality}
Pearl, J. (2009).
\newblock \emph{Causality: Models, Reasoning, and Inference}.
\newblock Cambridge University Press.

\bibitem{pearl2018book}
Pearl, J., \& Mackenzie, D. (2018).
\newblock \emph{The Book of Why}.
\newblock Basic Books.

\bibitem{marr1982vision}
Marr, D. (1982).
\newblock \emph{Vision}.
\newblock W. H. Freeman.

\end{thebibliography}

\end{document}
